\clearpage\pagestyle{otMain}\fontsize{11}{15}\selectfont\corrigesChapitre{8}{Fonctions : images et antécédents}
\begin{blocLivre}{cor-8-1}\begin{corrige}{1}{Calculer des images}\demarche{Substituer en gardant les parenthèses pour les négatifs.}
On remplace $x$ : $\res{f(0)=-4}$, $\res{f(2)=2}$ et $\res{f(-3)=-13}$. Chaque entrée donne une seule image.
\end{corrige}\end{blocLivre}
\begin{blocLivre}{cor-8-2}\begin{corrige}{2}{Lire un tableau}\demarche{Distinguer ce qui est connu dans le tableau et ce qui ne l’est pas.}
À la colonne de $-2$, l’image est $\res5$. La valeur 1 apparaît sous les entrées 0 et 3 : les antécédents indiqués sont $\res{0\text{ et }3}$. Le tableau ne dit pas s’il en existe ailleurs.
\end{corrige}\end{blocLivre}
\begin{blocLivre}{cor-8-3}\begin{corrige}{3}{Traduire une phrase}\demarche{Identifier le rôle de l’entrée et de la sortie.}
La première phrase s’écrit $\res{g(-1)=7}$. La deuxième signifie que 7 est l’image de 4, ou que 4 est un antécédent de 7. Les deux égalités peuvent être vraies pour une même fonction.
\end{corrige}\end{blocLivre}
\begin{blocLivre}{cor-8-4}\begin{corrige}{4}{Un point sur une courbe}\demarche{Calculer l’image de chaque abscisse.}
$f(3)=9-2=7$ : A appartient à la courbe. $f(-2)=4-2=2\ne1$ : B n’y appartient pas. Une lecture approximative ne remplace pas ces calculs.
\end{corrige}\end{blocLivre}
\begin{blocLivre}{cor-8-5}\begin{corrige}{5}{Un antécédent algébrique}\demarche{Écrire l’égalité correspondant à l’image demandée.}
On résout $5x+2=17$, d’où $5x=15$ et $\res{x=3}$. Vérification : $h(3)=17$. Le coefficient 5 non nul assure une seule solution.
\end{corrige}\end{blocLivre}
\begin{blocLivre}{cor-8-6}\begin{corrige}{6}{Une formule et un programme}\demarche{Traduire chaque étape avant de transformer l’expression.}
La fonction est $\res{f(x)=(x+3)^2}$. Ainsi $f(1)=16$ et $f(-3)=0$. Écrire $x^2+3$ n’aurait pas respecté l’ordre des instructions.
\end{corrige}\end{blocLivre}
\begin{blocLivre}{cor-8-7}\begin{corrige}{7}{Lire une ligne brisée}\demarche{Parcourir toute la courbe pour ne pas oublier une intersection.}
Le sommet d’abscisse $-1$ a pour ordonnée $\res3$. La droite horizontale d’ordonnée 1 rencontre la ligne aux abscisses $\res{-3,0,3}$. Sur le segment central, $f(x)=1-2x$, ce qui confirme $f(0)=1$.
\label{fig-cor-8-7}\figureGeo[corrige]{fonction-brisee}
\end{corrige}\end{blocLivre}
\begin{blocLivre}{cor-8-8}\begin{corrige}{8}{Plusieurs ou aucun antécédent}\demarche{Résoudre une équation de carré en examinant le signe.}
Les antécédents de 16 sont $\res{-4\text{ et }4}$. Celui de 0 est seulement $\res0$. Le nombre $-1$ n’a aucun antécédent réel, car un carré est positif ou nul.
\label{fig-cor-8-8}\figureGeo[corrige]{fonction-carre}
\end{corrige}\end{blocLivre}
\begin{blocLivre}{cor-8-9}\begin{corrige}{9}{Comparer deux représentations}\demarche{Relier une intersection à l’égalité des deux images.}
Le graphique suggère l’intersection $(3;5)$. L’équation $x+2=2x-1$ donne $\res{x=3}$ ; les deux images valent $\res5$. Le calcul valide la lecture.
\label{fig-cor-8-9}\figureGeo[corrige]{fonction-intersection}
\end{corrige}\end{blocLivre}
\begin{blocLivre}{cor-8-10}\begin{corrige}{10}{Ce que ne dit pas un tableau}\demarche{Repérer les informations absentes.}
Non dans les deux cas. L’entrée 2 n’est pas indiquée. Plusieurs courbes peuvent passer par les quatre points connus, avec des valeurs différentes en 2. Relier les points par des segments serait une hypothèse supplémentaire.
\end{corrige}\end{blocLivre}
\begin{blocLivre}{cor-8-11}\begin{corrige}{11}{Un réservoir}\demarche{Choisir les unités avant d’écrire la fonction.}
Le volume en m$^3$ vaut $\res{V(h)=0{,}8h}$, avec $0\le h\le1{,}5$. On convertit 600 L en 0,6 m$^3$. L’équation $0{,}8h=0{,}6$ donne $\res{h=0{,}75\text{ m}}$, qui respecte la hauteur maximale.
\label{fig-cor-8-11}\figureGeo[corrige]{cuve-hauteur}
\end{corrige}\end{blocLivre}
\begin{blocLivre}{cor-8-12}\begin{corrige}{12}{Un tarif plafonné}\demarche{Tenir compte du plafond et du domaine discret.}
Le prix est $2h$ tant qu’il ne dépasse pas 10 €, puis reste à 10 €. Les trois prix sont $\res{6\text{ €},10\text{ €},10\text{ €}}$. La fonction n’est pas linéaire sur tout l’ensemble : doubler 4 h en 8 h ne double pas le prix.
\end{corrige}\end{blocLivre}
\begin{blocLivre}{cor-8-13}\begin{corrige}{13}{Des aires possibles}\demarche{Exprimer l’autre dimension avec la contrainte de périmètre.}
L’autre côté vaut $20-x$. L’aire est $\res{A(x)=x(20-x)}$ cm$^2$, avec $0<x<20$ pour un rectangle non aplati. Les images de 3 et de 17 valent toutes deux $\res{51\text{ cm}^2}$. Une même aire peut donc avoir plusieurs antécédents.
\label{fig-cor-8-13}\figureGeo[corrige]{rectangle-fonction}
\end{corrige}\end{blocLivre}
\begin{blocLivre}{cor-8-14}\begin{corrige}{14}{Une lecture de déplacement}\demarche{Distinguer distance au départ et longueur totale du parcours.}
On s’éloigne jusqu’à 600 m pendant les dix premières minutes, on reste à cette distance cinq minutes, puis on revient en dix minutes. Dans le modèle rectiligne annoncé, la distance totale est $\res{1200\text{ m}}$. La dernière ordonnée 0 indique le retour au départ, pas un trajet total nul.
\label{fig-cor-8-14}\figureGeo[corrige]{promenade}
\end{corrige}\end{blocLivre}
\begin{blocLivre}{cor-8-15}\begin{corrige}{15}{Résoudre par un produit nul}\demarche{Choisir une forme adaptée à l’image recherchée.}
Le produit est nul pour $x-2=0$ ou $3x+6=0$, d’où $\res{2\text{ et }-2}$. Par exemple $f(0)=(-2)\times6=\res{-12}$, strictement négatif. Un exemple suffit pour cette dernière demande d’existence.
\end{corrige}\end{blocLivre}
